%------------------------------------------------------------------------------------
%  SECTION 6.3 : MISES EN PAGE MULTI-COLONNES (2 ET 3 COLONNES)
%------------------------------------------------------------------------------------
\section{Mises en Page Multi-Colonnes et Équilibrage Dynamique}
\label{sec:multi_colonnes}

Pour les documents techniques denses, les synthèses de laboratoire ou les articles de revue, une disposition en colonne unique (\emph{single-column}) peut engendrer des lignes trop longues (plus de 85 caractères), rendant la lecture fatigante.
Le passage en colonnes multiples permet de densifier l'information tout en conservant une longueur de ligne idéale (40 à 50 caractères).

\subsection{L'option globale \texttt{twocolumn} vs la flexibilité locale de \textsf{multicol}}
\label{subsec:comparatif_colonnes}

Il existe deux manières d'aborder le multi-colonnes en \LaTeX{} :
\begin{enumerate}
  \item \textbf{L'option globale de classe (\texttt{\textbackslash documentclass[twocolumn]\{ClasseGuillaume\}}) :}
  Elle scinde l'ensemble du mémoire en deux colonnes du début à la fin.
  Ce mode convient aux formats d'articles IEEE ou de comptes-rendus de conférences, mais s'avère trop rigide pour un mémoire universitaire dont la page de garde, les résumés et les grands tableaux doivent rester en pleine largeur.
  
  \item \textbf{L'approche locale dynamique avec le paquet \textsf{multicol} :}
  Elle permet de basculer ponctuellement en deux, trois ou quatre colonnes sur une fraction de page, puis de revenir à la colonne pleine de façon parfaitement fluide et sans saut de page forcé.
\end{enumerate}

\subsection{Démonstration en direct d'un passage en deux colonnes équilibrées}
\label{subsec:demo_deux_colonnes}

Voici une démonstration concrète d'un texte scindé en deux colonnes avec filet vertical séparateur délicat :

\setlength{\columnsep}{22pt}
\setlength{\columnseprule}{0.4pt}

\begin{multicols}{2}[\subsubsection*{Étude comparative des fluides caloporteurs en régime forcé}]
  Dans cette première colonne, nous examinons les propriétés thermophysiques de l'eau glycolée à \qty{30}{\percent}.
  Sa capacité thermique massique élevée ($c_p = \qty{3.85}{\kilo\joule\per\kilo\gram\per\kelvin}$) en fait un fluide de choix pour les circuits de refroidissement industriels soumis à des températures négatives modérées.
  
  Toutefois, sa viscosité dynamique augmente de manière sensible aux basses températures, entraînant une majoration des pertes de charge singulières dans les coudes de canalisations.
  
  \columnbreak % Forçage pédagogique de l'équilibrage de colonne
  
  Dans cette seconde colonne en regard, nous analysons le comportement d'une huile synthétique de synthèse.
  Bien que sa conductivité thermique soit inférieure ($\lambda = \qty{0.13}{\watt\per\meter\per\kelvin}$), sa très haute stabilité thermique jusqu'à \qty{320}{\degreeCelsius} autorise un fonctionnement sans pressurisation excessive du circuit primaire.
  
  L'équilibrage automatique assuré par le paquet \textsf{multicol} égalise précisément la hauteur des deux colonnes avant de restituer la largeur normale de page.
\end{multicols}

\setlength{\columnseprule}{0pt} % Rétablissement du séparateur invisible

\subsection{Démonstration d'un lexique technique en trois colonnes}
\label{subsec:demo_trois_colonnes}

Pour les nomenclatures, index de variables ou listes de constantes, le format en trois colonnes optimise l'occupation de la page :

\setlength{\columnsep}{16pt}
\begin{multicols}{3}
  \begin{itemize}[leftmargin=12pt]
    \item $\alpha$ : Diffusivité thermique
    \item $\beta$ : Dilatation volumique
    \item $\gamma$ : Indice adiabatique
    \item $\delta$ : Épaisseur de couche
    \item $\epsilon$ : Émissivité radiative
    \item $\zeta$ : Coefficient d'amortissement
    \item $\eta$ : Rendement isentropique
    \item $\theta$ : Température adimensionnelle
    \item $\kappa$ : Compressibilité isotherme
    \item $\lambda$ : Conductivité thermique
    \item $\mu$ : Viscosité dynamique
    \item $\nu$ : Viscosité cinématique
  \end{itemize}
\end{multicols}

